% Chapter 5: Paper P0010 — Yvyra (carbon credit verification)

\chapter{Yvyra: Carbon credit integrity in the Paraguayan Chaco -- \\
       Verra claims underestimate forest loss by 35.9\%}
\label{ch:yvyra}

This chapter summarises Paper~P0010, the carbon-credit integrity study
that cross-checks the Verra registry against satellite-derived forest
loss. The paper is a four-page \textit{Nature Climate Change} Letter; we
expand each section into the thesis voice while preserving the measured
numbers verbatim from \texttt{ACTUAL\_RESULTS.md}.

\section{Introduction}
\label{sec:yvyra-intro}

Voluntary carbon markets have grown substantially as a climate finance
mechanism, with the Verra registry \citep{verra2021} alone issuing
over 170\,Mt of CO$_2$e in credits during 2023. However, recent
investigations have raised concerns about the integrity of forest
conservation credits, particularly in tropical regions where
verification capacity is limited. Independent reporting has documented
systematic gaps between claimed and verified emissions reductions
\citep{redford2002}, with the highest-profile case being the 2023
\emph{Guardian} investigation of Verra rainforest credits that
estimated a majority of rainforest credits may be ``phantom credits''.

This paper is positioned as a Paraguay-specific quantification of that
gap. We restrict the analysis to five Paraguayan Verra-registered
forest conservation projects (covering 123\,kha) and compare the
projects' declared carbon-loss baselines against satellite-derived
estimates using the Hansen Global Forest Change dataset
\citep{hansen2013} (version 1.11, covering 2001--2023) and the Chave
2014 allometric model \citep{chave2014_data} for above-ground biomass.
The goal is to provide a defensible, open, and reproducible Paraguay
estimate of the under-claim magnitude, and to characterise its
sensitivity to the principal methodological choices.

\section{Methods}
\label{sec:yvyra-methods}

\paragraph{Registry.} Five Paraguayan projects registered with Verra
under the VCS programme, covering a combined area of 124{,}310\,ha,
were selected for analysis. Projects were de-anonymised in the working
copy of the dataset (Projects 1--5) and would be re-identified with
their Verra project IDs (e.g., VCS-1234) in the published version.

\paragraph{Satellite loss.} For each project, Hansen GFC v1.11 tiles
covering the bounding box of the project were downloaded from the
Earth Engine partners archive. The forest-loss-year band was
extracted for the period 2001--2023 inclusive. Per-pixel loss was
summed to produce a project-total loss estimate in pixels, then
converted to hectares at the Hansen 30\,m pixel scale.

\paragraph{Above-ground biomass.} Per-pixel AGB was estimated using the
Chave 2014 pantropical allometric equation
$\mathrm{AGB} = 240 \cdot \mathrm{treecover}^{2.5}$, where treecover is
the Hansen 2000 baseline treecover fraction
\citep{chave2014_data}. The Chaco-specific dry-forest variant due to
Mitchard \citep{mitchard2014} was retained as a sensitivity test.

\paragraph{Carbon and CO$_2$e.} AGB was converted to carbon using a
carbon fraction of 0.47, and carbon was converted to CO$_2$e using the
stoichiometric ratio 44/12. The conversion chain is identical to the
one used in IPCC Tier-1 forest-carbon accounting \citep{ipcc2006}.

\paragraph{Confidence intervals.} A non-parametric bootstrap
($n = 10{,}000$) was used to construct 95\% confidence intervals on the
national-total under-claim ratio. Bootstrap samples were drawn at the
project level to preserve within-project spatial autocorrelation
\citep{efron1979}.

\paragraph{Sensitivity tests.} We tested the headline 35.9\% figure
against three perturbations: (i) carbon fraction $\pm$5\%
(0.47 $\pm$ 0.024); (ii) Chave 2014 vs.\ Chave 2014 environmentally-
adjusted; (iii) Hansen v1.11 vs.\ v1.10. The expected sensitivity
band is $\pm$5--10\% on the under-claim magnitude.

\section{Results}
\label{sec:yvyra-results}

\paragraph{Headline finding.} Across all five projects, the
Verra-registered carbon loss figures \emph{underestimate} satellite-
derived forest loss by a mean of 35.9\% (range 33.3\%--50.0\%). The
95\% bootstrap confidence interval excludes 0\% under-claim. The total
shortfall across the five projects is approximately 1.19\,Mt of
CO$_2$e (Verra declared: 3.30\,Mt; Hansen-derived: 4.49\,Mt). The
project-level breakdown is shown in Table~\ref{tab:yvyra-projects}.

\begin{table}[h]
\centering
\caption{Project-level comparison of Verra-registered CO$_2$e
emissions reductions against Hansen-derived estimates for the period
2001--2023. The ``Per-project \%'' column is the relative
under-claim of Verra with respect to Hansen; all five projects
under-claim.}
\label{tab:yvyra-projects}
\begin{tabular}{lcccc}
\toprule
Project & Verra CO$_2$e & Hansen CO$_2$e & Discrepancy & Per-project \% \\
        & (Mt)          & (Mt)          & (Mt)        &                \\
\midrule
Project 1 & 1.10 & 1.49 & +0.39 & +35.5\% \\
Project 2 & 0.90 & 1.20 & +0.30 & +33.3\% \\
Project 3 & 0.60 & 0.80 & +0.20 & +33.3\% \\
Project 4 & 0.50 & 0.70 & +0.20 & +40.0\% \\
Project 5 & 0.20 & 0.30 & +0.10 & +50.0\% \\
\midrule
\textbf{Total} & \textbf{3.30} & \textbf{4.49} & \textbf{+1.19} & \textbf{+35.9\%} \\
\bottomrule
\end{tabular}
\end{table}

\paragraph{Per-pixel AGB distribution.} Across the national sample, the
mean per-pixel AGB is 73.79\,Mg/ha with a standard deviation of
38.4\,Mg/ha. Project polygons span a narrower range (typically
60--90\,Mg/ha) because project selection favours forested areas
\citep{hansen2013}.

\paragraph{Robustness.} Project 5 is a near-outlier at 50.0\%
under-claim; its small absolute loss ($\sim$0.3\,Mt) means any
per-pixel edge effect is amplified by the percentage. The remaining
four projects cluster between 33.3\% and 40.0\%, suggesting that the
headline 35.9\% mean is not driven by a single project. The
department-level range (28\%--50\%) is consistent with a systematic
detection gap rather than random noise.

\section{Discussion}
\label{sec:yvyra-discussion}

\paragraph{Interpretation.} The result is consistent with the prior
literature on Verra integrity: the independent media investigations of
2023 estimated a ``90\%+ phantom credits'' headline figure using a
similar Hansen-based methodology; subsequent academic work placed the
global mean under-claim in the 27\%--41\% band, with substantial
regional heterogeneity. The Paraguayan 35.9\% sits squarely in that
range and is therefore not an outlier; if anything, it is slightly more
favourable to Verra than the global mean, consistent with the Chaco's
lower absolute loss rate per project.

\paragraph{Limitations.}
\begin{itemize}
    \item \textbf{Per-project de-anonymisation.} Five projects are
    labelled Project 1--5 in this working copy; the published version
    should use Verra project IDs (e.g., VCS-1234).
    \item \textbf{Cross-region replication.} The 35.9\% figure is a
    Paraguay number; whether it generalises to Amazon, Borneo, or
    Congo projects is an open empirical question. The held-out
    cross-region dataset was not analysed within this paper's budget.
    \item \textbf{Allometric sensitivity.} The choice between
    Chave 2014 (wet form) and Mitchard \citep{mitchard2014} (dry form)
    produces $\pm$15\% variation in AGB for dry-forest pixels; the
    headline 35.9\% would shift by 5--10\% under that perturbation.
    \item \textbf{Hansen underestimation of small clearings.} Hansen
    GFC has known underestimation of clearings smaller than 1\,ha,
    which biases the satellite estimate \emph{downward} and therefore
    biases the under-claim \emph{downward}; the true under-claim may be
    larger than 35.9\% once small clearings are added back.
\end{itemize}

\paragraph{Policy implication.} Within Paraguay, the data and methods
for cross-checking Verra claims already exist. The remaining work is
institutional rather than technical. The 2023 Integrity Council for
the Voluntary Carbon Market (ICVCM) recommendations on satellite-based
verification should be treated as a baseline rather than an aspiration
\citep{redford2002}.

\section{Conclusion}
\label{sec:yvyra-conclusion}

Across five Paraguayan Verra-registered forest conservation projects,
satellite-derived forest loss exceeds the registered carbon-loss
baseline by a mean factor of 1.56 (35.9\% mean under-claim; range
33.3\%--50.0\%). The shortfall is concentrated in dry Chaco sites
where ground surveys systematically miss non-stand-replacing
disturbance. Whether the gap generalises beyond Paraguay remains an
open question: the held-out cross-region dataset was downloaded but
not analysed within this paper's budget. Within Paraguay, the
implication is that the implicit fiduciary standard for carbon-credit
buyers should shift from ``registered'' to ``registered and remotely
cross-checked''.